% 05-autograd.tex
\section{自动微分：显式有向无环图、就绪队列与版本守卫}
\label{sec:autograd}

autograd 引擎按代码逐条记录。没有 \ac{tape}：搜索该符号返回零。该图是通过 \texttt{Edge} 链接的 \texttt{shared\_ptr<Node>} 的显式 \ac{dag}，进行依赖项计数，并在由 \texttt{sequence\_nr} 键入的 ReadyQueues 之间进行调度。按顺序介绍节点、边、任务、输入缓冲区、就绪队列和引擎。

%==============================================================================
\subsection{无磁带，显式 DAG}
\label{subsec:notape}

引擎在转发过程中急切地构建图表。每个操作的自动微分包装器分配一个 \texttt{Node} 子类（从 via 生成到 \texttt{AutogradNodesGenerated.h} 的 148 种类型）或使用手写节点（例如\ \texttt{CatBackward}、\texttt{MeanBackward}、\texttt{AccumulateGrad}）。该节点的 \texttt{next\_edges:\ vector<Edge>} 指向其父节点。整个图是这些边的传递闭包。打印或绘制图表需要\texttt{shared\_ptr<Node>}；设备分区调度是从节点到队列的映射，无需扫描全局磁带。仅当 \texttt{roots.size!=1} 时才创建合成 \texttt{GraphRoot}，否则单根直接是图根。

%==============================================================================
\subsection{节点：字段和不变量}
\label{subsec:node}

是基础。 Listing~\ref{lst:nodefields} 提炼出布局；表~\ref{tab:nodefields} 解释了每个字段在热路径上的作用。

\begin{lstlisting}[caption={output metadata and base.}, label=lst:nodefields]
struct OutputSlotMeta {                   // 26
  std::vector<int64_t> shape;             // 27
  DType dtype{DType::Undefined};          // 28
  int64_t device_index = -1;              // 29
  bool valid = false;                     // 30
};
inline uint64_t get_and_increment_sequence_nr() { // 34
  static thread_local uint64_t counter = 0;       // 35
  return counter++;                               // 36
}
class TENSORPLAY_API Node : public std::enable_shared_from_this<Node> { // 44
  virtual variable_list apply(variable_list&&) = 0;  // 75
  void add_next_edge(Edge e){ update_topological_nr(e); next_edges_.push_back(std::move(e)); } //86
  void update_topological_nr(const Edge& e){         //141
    if(!e.is_valid()) return;
    auto topo=e.function->topological_nr();
    if(topological_nr_ <= topo) topological_nr_=topo+1; // 145
  }
protected:
  std::vector<Edge> next_edges_;               //150
  uint64_t sequence_nr_ = 0;                   //151 ctor assigns get_and_increment_sequence_nr()
  bool materialize_grads_ = true;              //152
  bool is_view_fn_ = false;                    //153
  bool multi_output_view_ = false;             //154
  std::string forward_op_name_;                //155
  std::vector<OutputSlotMeta> output_metas_;   //156
  uint64_t topological_nr_ = 0;                //157
private:
  std::vector<PreHookFn> pre_hooks_;           //160
  std::vector<PostHookFn> post_hooks_;         //161
  std::unique_ptr<AnomalyMetadata> anomaly_metadata_=nullptr; //162
};
\end{lstlisting}

\begin{table}[H]
\centering
\caption{\texttt{Node.h} 字段。除了冷的钩子/异常之外，每个字段都处于后向关键路径上。}
\label{tab:nodefields}
\small
\begin{tabular}{lll}
\toprule
场地 & Line & 目的 \\
\midrule
\texttt{next\_edges\_} & 150 & 父母。 \texttt{num\_inputs} 默认为 \texttt{size}。 \\
\texttt{sequence\_nr\_} & 151+50 & 线程本地单调 id 来自。订单 \texttt{ReadyQueue}。 \\
\texttt{topological\_nr\_} & 157+141 & 每个 \texttt{add\_next\_edge} 上的 \texttt{max(next.topo)+1}。启用 \texttt{min\_topo} 修剪。 \\
\texttt{materialize\_grads\_} & 152 & 如果为 false，则未定义的等级通过；如果为 true（默认），引擎将通过提示进行零填充。 \\
\texttt{output\_metas\_} & 156 & \texttt{PyNode} 的每个输出槽元数据（请参阅 \S\ref{subsec:pynode}）。生成的节点为空。 \\
\texttt{is\_view\_fn\_} & 153 & 查看包装器； \texttt{VariableTypeUtils.h} 中的 \texttt{check\_inplace} 拒绝叶视图突变。 \\
\texttt{multi\_output\_view\_} & 154 & \texttt{chunk} 视图；从不就地可变，通过 \texttt{forward\_op\_name\_} 报告。 \\
\texttt{post\_hooks\_} & 160 & \texttt{function<void(variable\_list\&\&)>} 列出；向后读取，而图是不可变的。 \\
\texttt{anomaly\_metadata\_} & 162 & 懒惰 \texttt{AnomalyMetadata};在节点创建时捕获堆栈。 \\
\bottomrule
\end{tabular}
\end{table}

关键不变量：

\begin{itemize}[leftmargin=1.2em]
\item \texttt{sequence\_nr\_} 在构造函数中从 a 分配。每个线程都是单调的，而不是全局的，因此最大堆近似于没有全局原子的反向创建顺序。比较一下。
\item \texttt{topological\_nr\_} 是距叶子的最长路径距离。扫描 \texttt{outputs[].function->topological\_nr}，并修剪 \texttt{topo < min\_topo}。所需子图之外的早期操作永远不会进入 \texttt{dependencies\_}。
\item \texttt{output\_metas\_} 存在是因为对于自定义函数，后向输入对应于前向 \emph{outputs}，而不是 \texttt{next\_edges} 输入。 in 为每个输出张量填充一个 \texttt{OutputSlotMeta} (shape/dtype/device\_index/valid)。生成的节点将其留空并依赖 \texttt{Edge} 提示。
\item 挂钩仅在图形构建时追加。引擎在 \texttt{evaluate\_function:426,446} 期间迭代它们，而没有其他线程改变节点，因此钩子本身不需要锁定（一旦 \texttt{execute} 启动，图在逻辑上是不可变的）。
\item \texttt{anomaly\_metadata\_} 被延迟分配并通过调用 \texttt{store\_stack} 进行初始化，如果设置了 \texttt{tls\_current\_evaluating\_node}，则调用 \texttt{assign\_parent}。这会将向后创建的节点的创建堆栈链接到其父节点。
\end{itemize}

根据保存的张量从存储 \texttt{SavedVariable} 成员生成节点；中的手写节点执行相同操作。每个节点的析构函数不必是虚拟的以进行调度； \texttt{CurrentEvalNodeGuard} 需要 \texttt{shared\_from\_this} 才能在异常模式下捕获父级。

%==============================================================================
\subsection{边缘：物化提示}
\label{subsec:edge}

不携带张量，仅提示向后可以 \texttt{ops::zeros} 丢失梯度而不触及前向张量（列表~\ref{lst:edge}）。

\begin{lstlisting}[caption={hints.}, label=lst:edge]
struct TENSORPLAY_API Edge{                    //14
  std::shared_ptr<Node> function; uint32_t input_nr; //15
  std::vector<int64_t> shape_hint; bool has_shape_hint=false; //21
  std::optional<DType> grad_dtype;                           //27
  std::optional<DeviceType> device_type_hint;               //33
  std::optional<int64_t> device_index_hint;                //34
  bool has_input_metadata() const {                         //36
    return has_shape_hint && grad_dtype && device_type_hint && device_index_hint; //37
  }
  Edge(std::shared_ptr<Node> f,uint32_t nr) :function(move(f)),input_nr(nr){} //42
  Edge(std::shared_ptr<Node> f,uint32_t nr,vector<int64_t> s):function(move(f)),input_nr(nr),shape_hint(move(s)),has_shape_hint(true){} //44
  Edge():function(nullptr),input_nr(0){}                     //49
  bool is_valid() const { return function!=nullptr; }        //51
};
\end{lstlisting}

\begin{table}[H]
\centering
\caption{边缘提示及其消费者在 \texttt{Engine::evaluate\_function} 中。}
\label{tab:edgehints}
\small
\begin{tabular}{lll}
\toprule
Hint & Line & 消费地点 \\
\midrule
\texttt{ has\_shape\_hint} & 21 & 实现 \texttt{ops::zeros(shape,\dots)} 和。通过 bool 区分 \texttt{} 和无提示。 \\
\texttt{grad\_dtype} & 27 & 不匹配毕业生的 Dtype 演员阵容；如果 \texttt{InferenceMode::is\_enabled} 则跳过。 \\
\texttt{device\_type\_hint} & 33 & 实现 \texttt{ops::zeros} 的设备；类型单独存储，因为 CPU DLpack 可能会报告。 \\
\texttt{device\_index\_hint} & 34 & 同一分配地点；合并成. \\
\bottomrule
\end{tabular}
\end{table}

语义：每次 Python 绑定或生成的包装器连接张量的 \texttt{grad\_fn} 时，它都会在来自使用者的 \texttt{Node} 的传出 \texttt{Edge} 上记录前向输入的 \texttt{device} 三元组。在向后期间，如果 \texttt{vars[i]} 未定义且 \texttt{func->materialize\_grads==true} ，则引擎更喜欢 \texttt{output\_metas\_[i]} （自定义函数）并回退到 \texttt{Edge} 三元组以分配 \texttt{ops::zeros} ，而无需进入 Python。

%==============================================================================
\subsection{GraphTask：每个向后共享状态}
\label{subsec:graphtask}

每次调用都会创建一个 \texttt{GraphTask}。多个工作线程可以同时评估同一任务的不同节点，因此一旦执行开始，所有共享簿记都由 \texttt{mutex\_} 保护（Listing~\ref{lst:graphtask}、Table~\ref{tab:graphtask}）。

\begin{lstlisting}[caption={per-backward context.}, label=lst:graphtask]
struct GraphTask{ //22
  struct ExecInfo{ //23
    struct Capture{ int input_idx_, output_idx_; Capture(int i,int o):input_idx_(i),output_idx_(o){} }; //24
    bool should_execute() const { return needed_ || captures_; } //30
    bool needed_=false; std::unique_ptr<vector<Capture>> captures_; //31
  };
  bool keep_graph_, grad_mode_; uint64_t trace_id_=0; //35+40
  unordered_map<Node*,InputBuffer> not_ready_;     //43
  unordered_map<Node*,int> dependencies_;           //44
  unordered_set<Node*> nodes_in_graph_;           //45
  unordered_map<Node*,ExecInfo> exec_info_;        //48 non-empty only for grad()
  vector<Tensor> captured_vars_;                    //50
  mutex mutex_; condition_variable cv_;            //53
  uint64_t outstanding_tasks_=0; bool completed_=false; exception_ptr exception_; //55
  explicit GraphTask(bool keep,bool grad):keep_graph_(keep),grad_mode_(grad){} //61
  void init_to_execute(Node& root, edge_list outputs, bool acc, uint64_t min_topo); //64
  void task_enqueued(){ lock_guard<mutex> l(mutex_); ++outstanding_tasks_; } //67
  bool task_completed(){ lock_guard<mutex> l(mutex_); if(--outstanding_tasks_==0){ completed_=true; cv_.notify_all(); return true; } return false; } //76
  void record_exception(exception_ptr e){ lock_guard<mutex> l(mutex_); if(!exception_) exception_=move(e);} //88
  bool is_completed(){ lock_guard<mutex> l(mutex_); return completed_; } //93
  void wait_for_completion(){ unique_lock<mutex> l(mutex_); cv_.wait(l,[this]{return completed_;}); } //99
};
\end{lstlisting}

\begin{table}[H]
\centering
\caption{\texttt{GraphTask.h} 字段；所有地图均由原始 \texttt{Node*} 键入（稳定，而 \texttt{shared\_ptr} 位于其他地方）。}
\label{tab:graphtask}
\small
\begin{tabular}{lll}
\toprule
场地 & Line & Role \\
\midrule
\texttt{keep\_graph\_} & 35 & 如果为 false，则调用 \texttt{func->release\_variables}（清除 \texttt{next\_edges\_}）。 \\
\texttt{grad\_mode\_} & 36 & \texttt{create\_graph} 决定；复制到每个工人的 \texttt{GradModeGuard} 中。 \\
\texttt{trace\_id\_} & 40 & 并发图的 \texttt{TP\_ENGINE\_TRACE} 过滤是单调的。 \\
\texttt{not\_ready\_} & 43 & 其 \texttt{dependencies>0} 的节点的缓冲区。由 \texttt{Node*} 键控，值 \texttt{InputBuffer}，大小为 \texttt{num\_inputs}。 \\
\texttt{dependencies\_} & 44 & 入度。计算一次，在 \texttt{mutex\_} in 下递减。 \\
\texttt{nodes\_in\_graph\_} & 45 & BFS 的重复数据删除集。 \\
\texttt{exec\_info\_} & 48 & 过滤 \texttt{grad}：空意味着 \texttt{backward} 执行所有内容。 \\
\texttt{captured\_vars\_} & 50 & \texttt{grad} 的输出：由引擎在 \texttt{apply} 之前从 \texttt{InputBuffer} 填充。 \\
\texttt{outstanding\_} & 53 & 共享状态保护、完成唤醒和任务计数器。 \texttt{outstanding\_tasks\_} 计数已入队但未完成的数量。 \\
\bottomrule
\end{tabular}
\end{table}

\subsubsection{使用 \texttt{init\_to\_execute} 进行过滤}
\label{subsubsec:inittoexec}

对于 \texttt{backward}，\texttt{exec\_info\_} 保持为空并且每个访问的节点都会执行。对于 \texttt{grad(outputs,inputs)}，请致电。该方法必须将 \emph{only} 标记为生成请求的输入所需的子图。

处的代码执行两次：

\begin{enumerate}[leftmargin=1.2em]
\item 当 \texttt{accumulate\_grad==true} 时，种子 \texttt{exec\_info\_[output\_edge.function].captures\_} 与 \texttt{(input\_nr, output\_idx)} 或 \texttt{needed\_=true}。预先确定结果向量的大小。
\item 使用 \texttt{Frame\{Node* fn; size\_t next\_next\_fn\_\}} 的显式堆栈从 \texttt{graph\_root} 开始遍历。一次迭代 \texttt{fn->next\_edges} 一个边。 \texttt{seen} 集加上 \texttt{nodeShouldExecute} lambda (\texttt{257}) 向后传播 \texttt{needed\_}：如果子级已经 \texttt{should\_execute}，则其父级被标记为 \texttt{needed\_=true} (\texttt{274})。当展开 (\texttt{287}) 时，刚刚完成的节点的 \texttt{needed\_} 被推送到其 DFS 父节点 (\texttt{289})。在（\texttt{outputs} 中的最小拓扑）处计算的 \texttt{min\_topo\_nr} 会修剪步行：任何具有 \texttt{topological\_nr < min\_topo} 的子级都会被跳过（\texttt{281}），因为它无法达到请求的输出。
\end{enumerate}

结果是，可以在复制捕获之后但在 \texttt{apply} 之前（当 \texttt{!fn\_info.needed\_} 时）提前返回：该节点的输出仍然需要作为捕获的输入，但节点本身不需要计算新的梯度。

%==============================================================================
\subsection{InputBuffer：累积和设备路由}
\label{subsec:inputbuffer}

既是累加器又是路由提示（Listing~\ref{lst:inputbuffer}）。

\begin{lstlisting}[caption={accumulation.}, label=lst:inputbuffer]
inline bool can_accumulate_inplace(const Tensor& v){ //16
  if(!v.is_contiguous()) return false;                           //17
  if(v.unsafeGetTensorImpl().use_count()!=2) return false;        //20
  if(!v.impl()->has_storage()) return false;                      //21
  if(v.impl()->storage().use_count()!=1) return false;            //22
  return true;                                                    //23
}
inline void accumulate(vector<Tensor>& buf,size_t pos,Tensor&& var,bool grad_mode){ //29
  auto& old=buf[pos];                                             //30
  if(grad_mode) buf[pos]=tpx::ops::add(old,var);                 //34 second-order
  else if(can_accumulate_inplace(old)) buf[pos]=old.add_(var); //36 in-place
  else buf[pos]=old+var;                                          //38 out-of-place
}
struct InputBuffer{ //43
  vector<Tensor> buffer;                                          //78
  void add(size_t pos,Tensor&& v,bool gm=false){                 //52
    if(pos>=buffer.size()||!v.defined()) return;                  //53
    if(!buffer[pos].defined()) buffer[pos]=move(v);               //55
    else accumulate(buffer,pos,move(v),gm);                       //57
  }
  int device_index() const{ //65
    for(auto& t:buffer) if(t.defined() && t.device().is_cuda()) return (int)t.device().index(); //67
    return -1; //71 CPU / unspecified
  }
  static variable_list variables(InputBuffer&& g){ return move(g.buffer); } //74
};
\end{lstlisting}

\begin{table}[H]
\centering
\caption{\texttt{can\_accumulate\_inplace} 三重条件。}
\label{tab:inplace}
\small
\begin{tabular}{ll}
\toprule
查看 & Why \\
\midrule
\texttt{is\_contiguous} & 无论如何，非连续的 \texttt{add\_} 都会触发副本；不合适的 \texttt{+} 更便宜。 \\
\texttt{unsafeGetTensorImpl.use\_count==2} & 来自 \texttt{impl} 的临时 \texttt{shared\_ptr} 副本算作 1； \texttt{==2} 表示调用者保存最后一个逻辑引用。 \\
\texttt{storage.use\_count==1} & 确保同一存储（视图）没有其他张量别名；变异会破坏别名。 \\
\bottomrule
\end{tabular}
\end{table}

是引擎使用的唯一累积入口点。位置 \texttt{pos} 的第一名毕业生搬入；第二届及以后的毕业生通过。在 \texttt{grad\_mode==true}（即\ \texttt{create\_graph}，在 \texttt{task.grad\_mode\_} 处传播）下，累积本身通过 \texttt{tpx::ops::add} 构建二阶图（重新进入调度并记录新的 \texttt{Node}）。否则，当 \texttt{can\_accumulate\_inplace} 成立时，引擎会尝试就地 \texttt{add\_}，否则回退到 \texttt{old+var} 进行分配。

扫描第一个 CUDA 张量。用它来选择目标 \texttt{ReadyQueue}: \texttt{dev>=0 ? queue\_for\_device(dev): \&cpu\_queue}。因此，放置是由第一个定义的输入的设备决定的，而不是由节点自己的设备决定的；这与梯度张量与其前向输入位于同一设备上的观察结果相匹配。

%==============================================================================
\subsection{ReadyQueue：优先级为 \texttt{sequence\_nr}}
\label{subsec:readyqueue}

实现调度程序队列（清单~\ref{lst:readyqueue}）。

\begin{lstlisting}[caption={priority queue.}, label=lst:readyqueue]
struct ReadyQueue{ //23
  struct NodeTask{ //25
    shared_ptr<Node> fn_; InputBuffer input_buffer_; GraphTask* graph_; //26
    NodeTask(shared_ptr<Node> f,InputBuffer b,GraphTask* g):fn_(move(f)),input_buffer_(move(b)),graph_(g){} //33
    bool operator<(const NodeTask& o) const { return fn_->sequence_nr() < o.fn_->sequence_nr(); } //37 max-heap
  };
  void push(NodeTask t){ lock_guard<mutex> l(mutex_); heap_.push(move(t)); cv_.notify_one(); } //42
  template<StopFn> optional<NodeTask> pop_until(const StopFn& stop){ //53
    unique_lock<mutex> lock(mutex_);
    for(;;){
      if(!heap_.empty()){ auto t=move(const_cast<NodeTask&>(heap_.top())); heap_.pop(); return t; } //58
      if(stop()) return nullopt; //62
      cv_.wait_for(lock, chrono::milliseconds(5)); //65 poll
    }
  }
  void notify(){ lock_guard<mutex> l(mutex_); cv_.notify_all(); } //74
  mutex mutex_; condition_variable cv_; priority_queue<NodeTask> heap_; //80
};
\end{lstlisting}

该队列是 \texttt{sequence\_nr\_} 的最大堆（\ac{dag} 执行顺序；请参阅下面的定义框）。由于节点是按正向顺序分配的，较大的 \texttt{sequence\_nr} 大致更接近输出，因此弹出最大值近似于反向拓扑顺序，而无需计算每个反向的拓扑排序。这是一种启发式的，而不是精确的：真正的拓扑顺序是通过 \texttt{dependencies\_} 计数强制执行的，这可以防止在父级完成之前执行。

\begin{tpdef}[Definition (Priority Heap vs.\ Topological Order)]
令 $G=(N,E)$ 为后向 DAG，$\sigma:N\rightarrow\mathbb{N}$ 将每个节点映射到其 \texttt{sequence\_nr} （前向分配顺序）。 \texttt{ReadyQueue} 弹出最大 $\sigma$ 的挂起节点，这等于反向转发顺序。执行顺序为 \emph{valid} ，当且仅当每个节点都在其所有父节点之后运行（按照 \texttt{dependencies\_} 入度计数）。堆实现有效性 \emph{without} 计算拓扑排序：节点仅在其最后一个父节点将其入队后才进入堆，因此任何弹出顺序通过构造都是有效的； $\sigma$ 只是将 \emph{within} 的有效集偏向反向-正向顺序，从而改善了保存的中间体的局部性。
\end{tpdef}

就绪队列有两种使用模式。设备工作线程调用 \texttt{pop\_until([]{return false;})} 并通过 5\,ms 轮询永远阻塞。发起线程调用 \texttt{cpu\_queue.pop\_until([\&graph]\{return graph.is\_completed;\})}，因此它可以在 \texttt{GraphTask} 完成后立即退出，即使 CPU 队列为空。 5\,ms 可以避免丢失唤醒（通知程序在 \texttt{notify} 之前获取互斥锁）并限制空闲延迟，而无需忙于旋转。通知程序在 \texttt{notify\_all} 之前获取锁，因此在 \texttt{pop\_until} 的空检查和等待之间通知不会丢失。

推入增量不会增加计数器本身；调用者必须已经将未完成任务计数增加到 \texttt{task.mutex\_} 下。这种顺序很重要：工作人员在完成时递减，启动程序在 \texttt{outstanding\_tasks\_==0} 上阻塞。

%==============================================================================
\subsection{引擎：队列、工作线程和跨设备路由}
\label{subsec:enginequeues}

拥有队列和工作线程。

\begin{lstlisting}[caption={engine.}, label=lst:engineh]
class TENSORPLAY_API Engine{ //85
  static Engine& get_default_engine(); //88 leaked singleton
  variable_list execute(edge_list roots,variable_list inputs,bool keep,bool create,bool acc,edge_list out); //95
  void compute_dependencies(Node* root,GraphTask& t,uint64_t min_topo); //105
  void evaluate_function(GraphTask& t,Node* func,InputBuffer& in,ReadyQueue& cpu,ReadyQueue* local); //111
  void enqueue_task(GraphTask& t,ReadyQueue::NodeTask&& nt,ReadyQueue& cpu,ReadyQueue* local); //116
  void execute_task(ReadyQueue::NodeTask&& t,ReadyQueue& cpu,ReadyQueue* local); //120
  ReadyQueue* queue_for_device(int idx); //122
  void worker_main(ReadyQueue& q); //124
  mutex queues_mutex_; //126
  unordered_map<int,ReadyQueue*> ready_queues_; //129 -1 CPU, >=0 CUDA
  unordered_map<int,thread> device_threads_; //130 leaked by design
  static thread_local int nested_depth_; //133
};
\end{lstlisting}

\begin{lstlisting}[caption={lazy queue and worker creation.}, label=lst:queuefordevice]
ReadyQueue* Engine::queue_for_device(int idx){ //151
  lock_guard<mutex> lock(queues_mutex_); //152
  auto it=ready_queues_.find(idx); if(it!=end) return it->second; //153
  auto* q=new ReadyQueue(); ready_queues_[idx]=q;          //156 leaked
  if(idx>=0) device_threads_.emplace(idx, thread([this,q]{ worker_main(*q); })); //158
  return q; //163
}
void Engine::worker_main(ReadyQueue& q){ //166
  for(;;){ auto task=q.pop_until([]{return false;}); if(!task) continue; execute_task(move(*task), q, nullptr); } //168
}
void Engine::enqueue_task(GraphTask& task, ReadyQueue::NodeTask&& nt, ReadyQueue& cpu, ReadyQueue* local){ //609
  task.task_enqueued(); //611
  if(local){ local->push(move(nt)); return; } //612 nested: stay on local
  int dev=nt.input_buffer_.device_index(); //616
  ReadyQueue* target=(dev>=0)? queue_for_device(dev) : &cpu; //617
  target->push(move(nt)); //618
}
\end{lstlisting}

设计注意事项：

\begin{itemize}[leftmargin=1.2em]
\item 故意泄漏单例（\texttt{new Engine} 从未释放），因为排队的任务和工作人员可能比静态析构函数寿命更长。队列本身也被泄露。
\item \texttt{ready\_queues\_} 将 \texttt{-1} 映射到 CPU 队列，将非负整数映射到 CUDA 设备索引。每个 CUDA 设备在第一次使用时都会生成一个持久工作线程。 CPU没有worker；启动器线程耗尽它。
\item 永远不会退出。它在 \texttt{pop\_until} 中阻塞并调用，捕获任何异常，将其记录在 \texttt{GraphTask::exception\_} via 中，并用 标记完成。
\item 是 \texttt{thread\_local}。当检测到 \texttt{nested=true} 时，它将每个任务路由到本地堆栈并耗尽它本身 (\texttt{684})。这使得可重入向后（向后的 \texttt{apply} 内的向后）无死锁：忙于外部图的工作人员无法阻塞内部图，并且内部图的任务永远不会泄漏到设备队列，工作人员将在设备队列中等待本身被阻塞的外部启动器。
\item 构建规则很重要：根目录和存储库根目录都是指向 \texttt{TP\_PACKAGE\_DIR} 的 CMake 二进制目录（请参阅 \texttt{CMakeLists.txt} \texttt{TP\_PACKAGE\_DIR}）。并发链接器损坏 \texttt{\_C.so}；请参阅 \texttt{README.md} 编译注释。
\end{itemize}

%==============================================================================
\subsection{建立依赖计数}
\label{subsec:dependencies}

\begin{lstlisting}[caption={BFS in-degree.}, label=lst:computedep]
void Engine::compute_dependencies(Node* root,GraphTask& task,uint64_t min_topo){ //192
  vector<Node*> queue{root}; auto& deps=task.dependencies_; auto& nodes=task.nodes_in_graph_; //194
  while(!queue.empty()){
    auto fn=queue.back(); queue.pop_back(); //198
    if(fn->topological_nr() < min_topo) continue; //201 prune
    for(auto& e: fn->next_edges()){ //204
      if(auto* next=e.function.get()){
        deps[next] += 1; //206
        if(was_inserted(nodes,next)) queue.push_back(next); //215
      }
    }
  }
}
\end{lstlisting}

这是从 \texttt{graph\_root} 开始的 DAG 上的 DFS（向量作为堆栈）。每条边都为子级存储的入度贡献 1。 \texttt{nodes\_in\_graph\_} 集确保每个节点扩展一次，但计数器仍然为每个传入边累加 1（即使该节点已经被看到，地图也会发生碰撞）。 \texttt{min\_topo} 防护 (\texttt{201}) 会跳过在第一个感兴趣的输出之前创建的整个子树；这对于 \texttt{grad} 至关重要，因为只有前向的后缀才重要。当 \texttt{TP\_ENGINE\_TRACE=2} 时，跟踪级别~2 在每个边缘发出。

是 \texttt{outputs[].function} 中最小的 \texttt{topological\_nr\_}。如果 \texttt{outputs} 为空（纯 \texttt{backward}），则返回 0 并且不会删除任何内容。

%==============================================================================
\subsection{编排向后：\texttt{Engine::execute}}
\label{subsec:execute}

Listing~\ref{lst:execute} 是完整的编排。

\begin{lstlisting}[caption={orchestration.}, label=lst:execute]
variable_list Engine::execute(const edge_list& roots,const variable_list& inputs, //621
  bool keep, bool create, bool acc, const edge_list& outputs){
  GraphTask task(keep,create); task.trace_id_=EngineTrace::next_graph_id(); //636
  bool skip=roots.size()==1; shared_ptr<Node> graph_root; //647
  if(skip) graph_root=roots[0].function; else graph_root=make_shared<GraphRoot>(roots,inputs); //648
  auto min_topo=compute_min_topological_nr(outputs); //655
  compute_dependencies(graph_root.get(), task, min_topo); //656
  if(!outputs.empty()) task.init_to_execute(*graph_root, outputs, acc, min_topo); //658
  const bool nested=nested_depth_>0; ReadyQueue local; ReadyQueue& cpu=*queue_for_device(-1); //662
  if(skip){ //669
    InputBuffer buf(roots[0].function->num_inputs()); //670
    buf.add(roots[0].input_nr, Tensor(inputs[0]), create); //671
    enqueue_task(task, NodeTask(roots[0].function, move(buf), &task), cpu, nested?&local:nullptr); //672
  } else enqueue_task(task, NodeTask(graph_root, InputBuffer(), &task), cpu, nested?&local:nullptr); //675
  DepthGuard depth(nested_depth_); //681
  if(nested){ while(!task.is_completed()){ auto t=local.pop_until([]{return false;}); if(!t) break; execute_task(move(*t), cpu, &local); } //683
  } else { while(!task.is_completed()){ auto t=cpu.pop_until([&task]{return task.is_completed();}); if(!t) break; execute_task(move(*t), cpu, nullptr); } task.wait_for_completion(); } //689
  if(task.exception_) rethrow_exception(task.exception_); //705
  return move(task.captured_vars_); //709
}
\end{lstlisting}

Steps:

\begin{enumerate}[leftmargin=1.2em]
\item 验证 \texttt{roots.size==inputs.size} (\texttt{629}) 以及 \texttt{grad} 具有非空输出 (\texttt{633})。
\item 从原子计数器创建并分配 \texttt{trace\_id\_}。
\item 如果需要扇入（向后多输出），则构建一个临时的 \texttt{GraphRoot} ；它保留根边缘并将种子梯度转发给正确的父级。
\item 计算依赖关系 (\S\ref{subsec:dependencies})，对于 \texttt{grad}，通过 \texttt{init\_to\_execute} 计算 \texttt{captures} 过滤器。
\item 将根入队。在单根快速路径中，种子梯度通过显式 \texttt{grad\_mode==create\_graph} 插入到 \texttt{input\_nr} 处，因此第一个 \texttt{InputBuffer::add} 已经知道是否构建二阶历史。
\item 分支在 \texttt{nested\_depth\_} 上。嵌套的backs使用本地队列；顶级启动程序本身会耗尽 CPU 队列，而 CUDA 工作线程会耗尽其设备队列。启动器也会在 CPU 队列耗尽后调用，因为某些工作可能仍在设备工作线程上进行。
\end{enumerate}

RAII：进入时为 \texttt{++depth}，退出时为 \texttt{--depth}，即使是通过异常也是如此，因此重入跟踪是异常安全的。

%==============================================================================
\subsection{核心循环：\texttt{evaluate\_function} 十个步骤}
\label{subsec:evaluate}

是每个节点的解释器。 Listing~\ref{lst:evaluatehead} 显示序言和尾声；这十个步骤的索引如下，并在小节中详细介绍。

\begin{lstlisting}[caption={prologue.}, label=lst:evaluatehead]
void Engine::evaluate_function(GraphTask& task, Node* func, InputBuffer& inputs, //333
  ReadyQueue& cpu, ReadyQueue* local){
  auto& exec_info=task.exec_info_; //335
  GradModeGuard grad_guard(task.grad_mode_); //339 Honour create_graph on workers
  if(!exec_info.empty()){ //341 grad() filtering
    auto& info=exec_info.at(func);
    if(auto* caps=info.captures_.get()){ lock_guard<mutex> L(task.mutex_); for(auto& c:*caps) task.captured_vars_[c.output_idx_]=inputs.buffer[c.input_idx_]; } //345
    if(!info.needed_){ if(task.task_completed()) queue_for_device(-1)->notify(); return; } //352
  }
  variable_list outputs; //361
  { //362 profiling + trace scope
    optional<prof::OpRecord> rec; if(prof::g_active.load(acquire)) rec.emplace(intern_name("backward::"+func->name())); //367
    // ... materialize, hooks, guard, apply, post hooks (see sub-steps)
  }
  // ... shape/dtype/nan + propagation (see below)
}
\end{lstlisting}

\subsubsection{步骤~1：\texttt{GradModeGuard}}
\label{subsubsec:gradguard}

\begin{lstlisting}[caption={honour \texttt{create\_graph} on workers.}, label=lst:gradguard]
struct GradModeGuard{ explicit GradModeGuard(bool e):prev_(GradMode::is_enabled()){ GradMode::set_enabled(e); } ~GradModeGuard(){ GradMode::set_enabled(prev_); } bool prev_; }; //21
\end{lstlisting}

\texttt{GradMode} 作为 \texttt{thread\_local bool enabled\_=true} 生活在 \texttt{p10} (, \S\ref{subsec:guards}) 中。工作线程继承操作系统赋予它们的任何 TLS，这与任务的 \texttt{create\_graph} 无关。在该节点评估期间，防护程序会使用 \texttt{task.grad\_mode\_} 覆盖工作线程的 TLS，因此嵌套累积和分析会看到正确的值。

\subsubsection{步骤~2：\texttt{ExecInfo} 捕获并提前退出}
\label{subsubsec:execinfo}

已在清单 ~\ref{lst:evaluatehead}:345 中显示。对于 \texttt{grad}，引擎将捕获的槽直接从输入缓冲区复制到 \texttt{task.mutex\_} 下的 \texttt{task.captured\_vars\_} 中，而不复制没有捕获的节点的整个缓冲区。如果 \texttt{!needed\_} 节点已完成：如果这是最后一个任务，则递减未完成的任务并唤醒发起者。

\subsubsection{步骤~3：具体化缺失的梯度}
\label{subsubsec:materialize}

\begin{lstlisting}[caption={missing grads.}, label=lst:materialize]
variable_list vars=InputBuffer::variables(move(inputs)); //384
if(func->materialize_grads()){ //392
  auto& metas=func->output_metas(); auto& edges=func->next_edges(); //393
  size_t n=min(vars.size(), metas.empty()?edges.size():metas.size()); //395
  for(size_t i=0;i<n;++i) if(!vars[i].defined()){ //398
    vector<int64_t> shape; DType dt=DType::Undefined; DeviceType dt_type=DeviceType::CPU; int64_t didx=-1; bool have=false; //399
    if(!metas.empty() && i<metas.size() && metas[i].valid){ shape=metas[i].shape; dt=metas[i].dtype; didx=metas[i].device_index; have=true; } //404
    else if(i<edges.size() && edges[i].has_shape_hint && edges[i].grad_dtype && edges[i].device_type_hint && edges[i].device_index_hint){ shape=edges[i].shape_hint; dt=*edges[i].grad_dtype; dt_type=*edges[i].device_type_hint; didx=*edges[i].device_index_hint; have=true; } //410
    if(!have) continue; //419
    Device dev(dt_type, dt_type==DeviceType::CPU?-1:didx); //421
    vars[i]=ops::zeros(shape,dt,dev); //423
  }
}
\end{lstlisting}

策略：类 \texttt{PyNode} 的自定义函数通过 \texttt{attach\_outputs} 在 \texttt{output\_metas\_} 中记录每个输出的元数据；生成的衍生节点将其留空并回退到每个输入 \texttt{Edge} 三元组。优先考虑 \texttt{output\_metas\_}，因为自定义函数的后向输入对应于前向输出（它们的 \texttt{num\_inputs} 等于 \texttt{output\_metas\_.size}，请参阅）。使用记录的 dtype 在记录的设备上分配，因此丢失的 grads 永远不需要 Python 交叉。

\subsubsection{步骤~4：预挂钩}
\label{subsubsec:prehooks}

\begin{lstlisting}[caption={pre-hooks.}, label=lst:prehooks]
for(auto& hook: func->pre_hooks()) vars=hook(move(vars)); //426
\end{lstlisting}

钩子是在图形构建时安装的 \texttt{std::function<variable\_list(variable\_list\&\&)>} 。 Python 绑定用 \texttt{py::gil\_scoped\_acquire} 包装它们，因此即使对于 CUDA 工作线程来说，钩子执行也是 GIL 安全的。

\subsubsection{步骤~5：在\texttt{CurrentEvalNodeGuard}下申请}
\label{subsubsec:apply}

\begin{lstlisting}[caption={apply.}, label=lst:apply]
CurrentEvalNodeGuard guard(AnomalyMode::is_enabled()? func->shared_from_this():nullptr); //432
if(AnomalyMode::is_enabled()){ //434
  try{ outputs=func->apply(move(vars)); }catch(...){ if(auto* m=func->anomaly_metadata()) m->print_stack(func->name()); throw; } //435
} else outputs=func->apply(move(vars)); //444
\end{lstlisting}

将 \texttt{func} 推到 \texttt{tls\_current\_evaluating\_node} 上，因此在此 \texttt{apply} 内分配的任何 \texttt{Node} （可能在 \texttt{create\_graph} 下）将当前节点记录为其异常父节点。在启用异常的异常情况下，引擎会在重新抛出之前打印前向堆栈链。

还发出具有渲染形状/dtype 的 \texttt{"apply \%s inputs=[...]"} ，并且在级别~3 处最多八个元素值。

\subsubsection{步骤~6：后钩}
\label{subsubsec:posthooks}

\begin{lstlisting}[caption={post-hooks.}, label=lst:posthooks]
for(auto& hook: func->post_hooks()) outputs=hook(outputs, move(outputs)); //447 -- signature (inputs,outputs)->outputs
\end{lstlisting}

Post-hooks 接收该节点的输入和输出，并可能重写输出（例如\ grad Clipping）。 C++ 绑定公开 \texttt{(inputs,outputs)->outputs} 签名以及 GIL 获取。

\subsubsection{步骤~7：形状验证和\texttt{sum\_to\_shape}}
\label{subsubsec:sumto}

转发期间的广播将标量梯度膨胀为大形状。在下一个边缘消耗梯度之前，引擎必须将其减少到该边缘上记录的前向输入形状。列出~\ref{lst:sumtoshape} 是例行公事。

\begin{lstlisting}[caption={broadcast recovery.}, label=lst:sumtoshape]
static Tensor sum_to_shape(const Tensor& grad,const vector<int64_t>& target){ //298
  if(target.empty()) return ops::sum(grad); //301 scalar ()
  int64_t leading=(int64_t)grad.dim()-(int64_t)target.size(); //304
  if(leading<0){ int64_t n=1; for(auto d:target) n*=d; if(grad.numel()==n) return ops::reshape(grad,target); return grad; } //306
  vector<int64_t> reduce; for(int64_t i=0;i<leading;++i) reduce.push_back(i); //318
  for(int64_t i=leading;i<(int64_t)grad.dim();++i) if(target[i-leading]==1 && grad.size(i)!=1) reduce.push_back(i); //319
  Tensor cur=reduce.empty()?grad:ops::sum(grad,reduce,true); //325 keepdim
  if(leading>0) cur=ops::reshape(cur,target); //327
  return cur; //330
}
\end{lstlisting}

为每个 \texttt{i < min(outputs.size,edges.size)} 调用，如果 \texttt{edges[i].has\_shape\_hint} 和毕业生的尺寸与提示不匹配，则调用 \texttt{sum\_to\_shape}。该逻辑将维度保留在 \texttt{target==1} 但 \texttt{grad>1} 处，并与 \texttt{keepdim=true} 求和，然后重塑前导广播维度。在 \texttt{create\_graph} 下，\texttt{reshape} 调用连接二阶图，因为 \texttt{GradModeGuard} 仍然处于活动状态。

\subsubsection{步骤~8：Dtype Cast}
\label{subsubsec:dtypecast}

\begin{lstlisting}[caption={dtype alignment.}, label=lst:dtypecast]
auto& gdt=edges[i].grad_dtype;
if(gdt && isFloatingType(*gdt) && isFloatingType(outputs[i].dtype()) && outputs[i].dtype()!=*gdt && !InferenceMode::is_enabled())
  outputs[i]=tpx::to(outputs[i], *gdt); //501 autocast rebalance: fp32->bf16 etc.
\end{lstlisting}

存在边缘 dtype 提示，以便由未包装的提升操作生成的 \texttt{fp32} grad 重新进入其保存的张量是低精度的向后节点（自动转换情况记录在）。非浮动提示（布尔掩码、索引）永远不会是梯度，并且会被单独保留。 \texttt{InferenceMode} 防护与调度程序的 \texttt{inference\_tensor\_} 语义相匹配。

\subsubsection{步骤~9：NaN 探测}
\label{subsubsec:nanprobe}

\begin{lstlisting}[caption={NaN probe.}, label=lst:nanprobe]
bool tensor_has_nan(const Tensor& t){ //139
  if(!t.defined()|| !isFloatingOrComplexType(t.dtype())) return false; //140
  return t.ne(t).any().item<bool>(); //141 ieee NaN != NaN via dispatch
}
\end{lstlisting}

At：如果 \texttt{AnomalyMode::is\_enabled\&\&should\_check\_nan}，则在本地 \texttt{GradModeGuard(false)} 下探测每个输出，因此探测器本身不记录历史记录。发生故障时，发动机会抛出故障。

\subsubsection{步骤~10：传播、累加和入队}
\label{subsubsec:propagate}

列出 ~\ref{lst:propagate} 将输出分发给父母。

\begin{lstlisting}[caption={propagation.}, label=lst:propagate]
auto num=outputs.size(); auto& edges=func->next_edges(); //515
for(size_t i=0;i<num;++i){ //519
  if(i>=edges.size()) break; auto& out=outputs[i]; auto& nxt=edges[i]; if(!nxt.is_valid()) continue; //520
  bool ready=false, enqueue=false; ReadyQueue::NodeTask pending(nullptr,InputBuffer(),&task); //529
  { lock_guard<mutex> L(task.mutex_); //532
    auto it=task.dependencies_.find(nxt.function.get());
    if(it==end) TP_THROW(RuntimeError,"dependency not found"); else if(--it->second==0){ task.dependencies_.erase(it); ready=true; } //535
    auto& not_ready=task.not_ready_;
    auto jt=not_ready.find(nxt.function.get());
    if(jt==end){
      if(!exec_info.empty()){ auto k=exec_info.find(nxt.function.get()); if(k==end||!k->second.should_execute()) continue; } //546
      InputBuffer buf(nxt.function->num_inputs()); buf.add(nxt.input_nr, move(out), task.grad_mode_); //554
      if(ready){ pending=NodeTask(nxt.function, move(buf), &task); enqueue=true; } else not_ready.emplace(nxt.function.get(), move(buf)); //556
    } else {
      auto& buf=jt->second; buf.add(nxt.input_nr, move(out), task.grad_mode_); //564
      if(ready){ pending=NodeTask(nxt.function, move(buf), &task); enqueue=true; not_ready.erase(jt); } //566
    }
  }
  if(enqueue) enqueue_task(task, move(pending), cpu, local); //583
}
if(!task.keep_graph_) func->release_variables(); //594 clears next_edges_
if(task.task_completed()) queue_for_device(-1)->notify(); //598 wakes initiator
\end{lstlisting}

在 \texttt{task.mutex\_} 下，引擎递减 \texttt{dependencies\_[next]}，将第一个贡献存储在 \texttt{not\_ready\_[next]} 中，并为后续贡献调用 \texttt{InputBuffer::add:554,565}。当计数器为零时，缓冲区将移入 \texttt{NodeTask} 并交给 \texttt{enqueue\_task}，后者按 \texttt{device\_index} 选择队列（或嵌套图的本地队列）。清除 \texttt{next\_edges\_}，以便在 \texttt{retain\_graph==false} 时立即释放 DAG，但在 \texttt{create\_graph} 或 \texttt{retain\_graph} 请求时保持完整。

%==============================================================================
\subsection{SavedVariable 和 VariableVersion：突变防护}
\label{subsec:savedvar}

后向需要前向张量。如果张量在向前和向后之间就地变异，则存储裸露的 \texttt{Tensor} 会默默地产生错误的梯度。 TensorPlay 在存储上使用版本计数器以及受保护的句柄（清单~\ref{lst:savedvar}--\ref{lst:varver}）。

\begin{lstlisting}[caption={guarded handle.}, label=lst:savedvar]
class SavedVariable{ //15
  Tensor data_; uint32_t saved_version_=0; //39
  void save(const Tensor& t){ //8 in SavedVariable.cpp
    if(!t.defined()){ data_=Tensor(); saved_version_=0; return; } //9
    data_=t; saved_version_=t.unsafeGetTensorImpl()->version(); //14
  }
  Tensor unpack() const{ //18
    if(!data_.defined()) return Tensor(); //19
    uint32_t cur=data_.unsafeGetTensorImpl()->version(); //20
    if(cur!=saved_version_) TP_THROW(RuntimeError, //21
      "one of the variables needed for gradient computation has been modified by an inplace operation: [saved version: "+to_string(saved_version_)+"; current version: "+to_string(cur)+"]"); //22
    return data_; //28
  }
  void reset_data(){ data_=Tensor(); saved_version_=0; } //31 called by Node::release_variables
  bool defined() const { return data_.defined(); } //36
};
\end{lstlisting}

\begin{lstlisting}[caption={atomic counter.}, label=lst:varver]
struct VersionCounter{ atomic<uint32_t> version_{0}; void bump() noexcept { version_.fetch_add(1, relaxed); } uint32_t current_version() const noexcept { return version_.load(relaxed); } }; //13
class VariableVersion{ //24
  shared_ptr<VersionCounter> counter_; bool enabled_=true; //26
  VariableVersion():counter_(make_shared<VersionCounter>()){} //30 eagerly allocated
  explicit VariableVersion(bool e):counter_(e?make_shared<VersionCounter>():nullptr),enabled_(e){} //32
  uint32_t current_version() const { return counter_?counter_->current_version():0; } //36
  void bump(){ if(!enabled_) return; counter_->bump(); } //44
  bool shares_counter(const VariableVersion& o) const { return counter_ && counter_==o.counter_; } //56 alias test
};
\end{lstlisting}

\begin{table}[H]
\centering
\caption{版本生命周期。每条路径：线都是合同，而不是建议。}
\label{tab:verlife}
\small
\begin{tabular}{lll}
\toprule
Site & 路径：线 & 行动 \\
\midrule
\texttt{TensorImpl::version\_counter\_} &  & 初始化为\texttt{VariableVersion(!InferenceMode::is\_enabled)}；推理张量得到一个禁用的计数器。 \\
\texttt{TensorImpl::bump\_version} &  & \texttt{version\_counter\_.bump}（+ 推理模式防护）。 \\
生成的变异操作 &  & \texttt{if(!InferenceMode::is\_enabled) \_\_tp\_tensor.impl->bump\_version} 对于每个 \texttt{a!} 参数。 \\
查看创建 &  & \texttt{result.impl->share\_version\_counter(input.impl)} 所以视图+基共享一个 \texttt{VersionCounter}。 \\
\texttt{SavedVariable::save} &  & 快照\texttt{tensor.unsafeGetTensorImpl->version}。 \\
\texttt{SavedVariable::unpack} &  & 比较当前版本；如果消息中的两个版本不匹配，则会抛出异常。 \\
  &  & \texttt{is\_view\_of\_leaf} 在遍历视图节点之前检查 \texttt{attr\_version}。 \\
\bottomrule
\end{tabular}
\end{table}

为什么急切分配：总是分配 \texttt{VersionCounter}。如果分配是惰性的，则尚未突变的基数和在任何突变之前创建的视图将在第一次碰撞时各自增加自己的计数器，并且一个别名上的后续突变将对另一个别名不可见。热切分配加上保证所有别名观察到相同的凹凸。计数器是原子的，因此持有另一个别名的读者可以看到来自一个别名的碰撞。

%==============================================================================
\subsection{异常和模式防护}
\label{subsec:guards}

及其同伴将全局开关与线程本地当前节点分开（Listing~\ref{lst:anomaly}）。

\begin{lstlisting}[caption={and .}, label=lst:anomaly]
class AnomalyMode{ //14
  static bool is_enabled(){ return _enabled; } //16
  static bool should_check_nan(){ return _check_nan; } //17
  static void set_enabled(bool e,bool c=true){ _enabled=e; _check_nan=c; } //18
  static bool _enabled, _check_nan; //24 global, not TLS
};
struct AnomalyMetadata{ //32
  void store_stack(); void print_stack(const string& node); //39
  void assign_parent(const shared_ptr<Node>& p){ parent_=p; } //42 weak_ptr<Node> parent_;
  string traceback_; weak_ptr<Node> parent_; //45
};
namespace{ thread_local weak_ptr<Node> tls_current_evaluating_node; } //20 in .cpp
shared_ptr<Node> get_current_evaluating_node(){ return tls_current_evaluating_node.lock(); } //28
void set_current_evaluating_node(const shared_ptr<Node>& n){ tls_current_evaluating_node=n; } //32
struct CurrentEvalNodeGuard{ //57
  explicit CurrentEvalNodeGuard(const shared_ptr<Node>& n):last_(get_current_evaluating_node()){ set_current_evaluating_node(n); } //59
  ~CurrentEvalNodeGuard(){ set_current_evaluating_node(last_); } //62
  shared_ptr<Node> last_; //69
};
\end{lstlisting}

状态机：

\begin{itemize}[leftmargin=1.2em]
\item 用户调用 \texttt{tensorplay.autograd.set\_detect\_anomaly(True)} 到达 \texttt{AnomalyMode::set\_enabled(enabled,check\_nan)}。这设置了两个对每个线程可见的全局变量；没有 TLS，没有引用计数，因此第一个启用它的线程将支付所有费用。
\item 创建节点时，通过 \texttt{store\_stack:38} 快照当前堆栈。默认情况下调用 \texttt{tensorplay::get\_stacktrace} （C++ 框架）。 Python 绑定安装了一个在 GIL 下调用 \texttt{traceback.format\_stack} 的捕获器，因此存储的回溯是前向操作的 Python 调用站点，而不是 C++ 回溯。
\item 在向后期间，将评估节点推到 \texttt{tls\_current\_evaluating\_node} 上。在该 \texttt{apply}（二阶图）内创建的任何节点都会将推送的节点记录为其 \texttt{AnomalyMetadata::parent\_}。失败时，引擎会在 \texttt{print\_stack} 处遍历链，为每个父项打印 \texttt{"Previous calculation was induced by..."} ，直到弱指针过期。
\item NaN 检查（通过 \texttt{t.ne(t).any}）被门控并在本地 \texttt{GradModeGuard(false)} 下运行，因此探针永远不会改变 \texttt{GradMode}。
\end{itemize}

\subsubsection{GradMode 和 InferenceMode：线程本地开关}
\label{subsubsec:gradmode}

\begin{lstlisting}[caption={and .}, label=lst:gradmode]
class GradMode{ static bool is_enabled(){ return enabled_; } static void set_enabled(bool e){ enabled_=e; } static thread_local bool enabled_; }; //8 in GradMode.h
// GradMode.cpp:5  thread_local bool GradMode::enabled_ = true; // enabled by default
class InferenceMode{ static bool is_enabled(){ return enabled_; } static void set_enabled(bool e){ enabled_=e; } static thread_local bool enabled_; }; //9
// InferenceMode.cpp:5  thread_local bool InferenceMode::enabled_ = false;
\end{lstlisting}

两者都故意存在于 \texttt{p10} 中（Python 绑定将它们重新导出为 \texttt{ InferenceMode}），因此调度可以在不依赖 \texttt{tpx} 的情况下查阅它们。它们是真实的 \texttt{thread\_local}: 一个线程中的 \texttt{no\_grad} 不会泄漏到引擎工作线程或其他用户线程中。推理张量天生具有 \texttt{version\_counter\_\{false\}} 和 \texttt{inference\_tensor\_=true}；省略它们的 Autograd 位，并在 \texttt{InferenceMode::is\_enabled} 时跳过 dtype 转换。 RAII 帮助程序并恢复以前的值，包括异常情况。

%==============================================================================
\subsection{Python 绑定：\texttt{PyNode} 和单交叉热路径}
\label{subsec:pynode}

自定义 autograd 函数是用户代码的逃生口。绑定位于 (Listing~\ref{lst:pynode})。

\begin{lstlisting}[caption={\texttt{PyNode}.}, label=lst:pynode]
class PyNode : public tpx::Node{ //24
  PyNode(py::object ctx): py_ctx_(move(ctx)){} //26
  size_t num_inputs() const override { return output_metas().empty()? Node::num_inputs() : output_metas().size(); } //30 Py outputs are backward inputs
  variable_list apply(variable_list&& inputs) override{ //35
    py::gil_scoped_acquire gil; //37 workers need the GIL
    py::tuple py_in(inputs.size()); for(size_t i=0;i<inputs.size();++i) py_in[i]=inputs[i].defined()? py::cast(inputs[i]): py::none(); //42
    py::object res=py_ctx_.attr("backward")(*py_in); //58
    variable_list out; if(res.is_none()) return out; else if(py::isinstance<Tensor>(res)) out.push_back(py::cast<Tensor>(res)); else for(auto item: py::cast<py::sequence>(res)) out.push_back(item.is_none()? Tensor(): py::cast<Tensor>(item)); //60
    return out; //79
  }
  py::object py_ctx_; //82
};
\end{lstlisting}

\begin{lstlisting}[caption={\texttt{attach\_outputs} records hints for zero-fill.}, label=lst:attach]
.def("attach_outputs", [](PyNode& self, py::handle outputs){ //159
  auto node=static_pointer_cast<Node>(self.shared_from_this()); //166
  auto& metas=self.output_metas(); metas.clear(); auto record=[&](py::handle item){ //169
    OutputSlotMeta m; if(py::isinstance<Tensor>(item)){ auto& t=py::cast<const Tensor&>(item); m.shape=vector<int64_t>(t.shape()); m.dtype=t.dtype(); m.device_index=t.device().index(); m.valid=true; } metas.push_back(move(m)); //173
  };
  if(py::isinstance<Tensor>(outputs)){ record(outputs); Tensor& t=py::cast<Tensor&>(outputs); impl::set_requires_grad(t,true); impl::set_grad_fn(t,node,0); return; } //183
  for(auto item: outputs.cast<py::sequence>()){ record(item); if(py::isinstance<Tensor>(item)){ Tensor& t=py::cast<Tensor&>(item); impl::set_requires_grad(t,true); impl::set_grad_fn(t,node,idx); } ++idx; } //190
}, "outputs"_a);
\end{lstlisting}

绑定不仅仅是一个蹦床。三个热路径优化减少了每个操作的 Python 开销：

\begin{enumerate}[leftmargin=1.2em]
\item 和批量边缘连接：一次交叉通过迭代 \texttt{py::isinstance<Tensor>} 并调用每个张量，然后 \texttt{add\_next\_edge\_list} 一次，为所有张量参数构建 \texttt{vector<Edge>} 。否则，非批处理的 Python 循环将付出 N 次 pybind 交叉。
\item 是单入口快速路径。它创建 \texttt{PyNode}、解压 \texttt{needs\_input\_grad}、连线 \texttt{next\_edges} (\texttt{327})、禁用 GradMode (\texttt{346})、调用 \texttt{forward} 加 \texttt{setup\_context}、用 \texttt{set\_grad\_fn} 加 \texttt{OutputSlotMeta} (\texttt{368}) 标记输出，并返回 \texttt{(output, ctx, needs, executable, node)}，所有这些都在一个 GIL 保持中完成。 Python 之后只构建向后闭包。这将五次往返替换为一次。
\item \texttt{grad} 在进入 \texttt{grad} 之前释放 GIL，因为工作人员需要它才能重新进入 Python \texttt{PyNode::apply}。坚持下去就会陷入僵局。
\end{enumerate}

当 \texttt{output\_metas\_} 非空时覆盖默认的 \texttt{next\_edges\_.size} ，因为对于自定义函数，后向元数是前向的输出元数，而不是其输入元数。 \texttt{Node} 上的钩子注册使用 GIL 获取包装用户钩子，因此即使从 CUDA 工作线程调用钩子也是安全的。

%==============================================================================
\subsection{深入探讨：调度故事下的六种机制}
\label{subsec:autograd-deep}

前面的部分描述了引擎的数据结构和十步解释器。六个更精细的机制——代码中的每一个小，实践中的每一个承载——完成了调度程序如何在修剪、叶子、失败、钩子、重入和观察下保持正确的图景。

%----------------------------------------------------------------------
\subsubsection{\texttt{topological\_nr} 修剪栅栏}
\label{subsubsec:topoprune}

除了 \texttt{sequence\_nr} 之外，每个节点还携带第二个计数器：在边缘插入时增量维护的 \texttt{topological\_nr} 。当一个节点链接到父节点时，它的计数器变成父节点的计数器加一——如果它大于它的当前值：

\begin{lstlisting}[caption={Topological number maintenance is two lines at graph-construction time---no pass over the graph is ever needed later.}]
void update_topological_nr(const Edge& edge) {
    if (!edge.is_valid()) return;
    auto topo_nr = edge.function->topological_nr();
    if (topological_nr_ <= topo_nr) {
        topological_nr_ = topo_nr + 1;
    }
}
\end{lstlisting}

计数器使引擎能够进行最有效的优化：在对依赖关系进行计数之前，引擎会计算所请求的输出中的 \emph{minimum} 拓扑编号，并且 BFS 和执行标记遍历都会跳过编号较小的任何节点。发出这种声音的不变量是： \texttt{topo\_nr} 小于输出的节点可以通过构造，永远不会位于输出的路径 \emph{to} 上，因此步行不需要访问它。在 \texttt{backward()} 仅针对后期输出的图表上，早期前向操作的整个子树在计数器构建之前就被修剪——栅栏是阻止 \texttt{grad(outputs=[y])} 在整个图表上为 \texttt{backward()} 付费的原因。

\subsubsection{\texttt{AccumulateGrad}：叶子合约}
\label{subsubsec:accumulategrad}

叶张量不通过通用传播路径接收梯度。每个带有 \texttt{requires\_grad} 的叶子都拥有一个专用的 \texttt{AccumulateGrad} 节点，其 \texttt{sequence\_nr} 是 \texttt{UINT64\_MAX}，强制优先级队列在其依赖项允许的情况下尽早对其进行调度。它的主体在一页中实现了完整的叶契约：设备检查、形状协调、连续性和最终存储：

\begin{lstlisting}[caption={The leaf node's shape reconciliation: the same reduce-then-reshape logic as the engine-wide \texttt{sum\_to\_shape}, specialized for parameters.}]
explicit AccumulateGrad(const Tensor& t)
    : Node(UINT64_MAX), value_(t) {}   // always scheduled ASAP

variable_list apply(variable_list&& inputs) override {
    if (inputs.empty() || !inputs[0].defined()) return {};
    Tensor grad = inputs[0];
    /* device check: param vs grad must agree */
    /* shape reconciliation:
       1. sum away extra leading dims (broadcast inflated them)
       2. sum away dims where param has size 1 and grad larger
       3. reshape to the parameter's exact rank and sizes      */
    if (auto* meta = impl::get_autograd_meta(value_)) {
        // Materialize strided gradients: downstream consumers (foreach
        // optimizers, .numpy()) rely on dense storage.
        if (!grad.is_contiguous()) grad = grad.contiguous();
        meta->accum_grad(grad);
    }
    return {};
}
\end{lstlisting}

三个设计要点值得一读。首先，形状修复存在于 \emph{at the leaf} 中，而不是在每个生产节点中：中间节点沿其边缘传递广播膨胀的形状，并且只有最终目的地进行协调。其次，连续性物化记录了它的消费者契约——\emph{foreach}优化器内核和NumPy桥都需要密集存储，因此转置梯度在这里复制一次，而不是稍后让消费者感到惊讶。第三，存储的元数据端是弱引用握手：张量的 \texttt{grad\_accumulator} 是 \texttt{weak\_ptr<Node>} ，正是因为节点强烈拥有该张量——强边回将形成一个无法收集的循环，泄漏参与图中的每个叶子。

\subsubsection{异常遏制和首次错误规则}
\label{subsubsec:exceptions}

向后是多个工作线程上的并发执行，因此异常处理是一个设计问题，而不是事后的想法。规则：（i）每个工作线程捕获节点的 \texttt{apply} 抛出的每个异常，并以先胜策略将其记录到 \texttt{GraphTask} 中——互斥锁保护存储保留 \emph{first} 异常指针并默默地删除后面的异常指针，因为已经失败的图不需要三个错误报告； (ii) 记录工作人员然后精确地减少未完成任务计数，就像成功时一样，因此完成计数在失败路径上保持正确； (iii) 启动线程阻塞任务的条件变量，在完成时唤醒，并在图完全停顿后重新抛出记录的异常。结果是：一个节点的形状错误永远不会让同级工作线程旋转，并且在图耗尽之后，用户会看到来自单个确定性位置（启动线程）的单个回溯。

\begin{lstlisting}[caption={Completion accounting is identical for success and failure; only the initiating thread ever rethrows.}]
bool task_completed() {
    std::lock_guard<std::mutex> lock(mutex_);
    --outstanding_tasks_;
    if (outstanding_tasks_ == 0) {
        completed_ = true;
        cv_.notify_all();       // wakes the initiator's wait_for_completion
        return true;
    }
    return false;
}
void record_exception(std::exception_ptr exc) {
    std::lock_guard<std::mutex> lock(mutex_);
    if (!exception_) exception_ = std::move(exc);   // first error wins
}
\end{lstlisting}

\subsubsection{最终回调：没有竞赛的后向后钩子}
\label{subsubsec:finalcallbacks}

一些消费者需要代码来运行 \emph{after} 整个图完成，但 \emph{from inside} 向后调用——FSDP 向后完成挂钩就是一个激励性的例子。普通的“在 \texttt{backward} 末尾调用我的函数”API 无法在任何工作人员可能最后完成的引擎中工作。引擎公开 \texttt{queue\_callback}，仅在当前线程上执行向后操作时有效； \texttt{GraphTask} 将回调向量存储在其自己的互斥体下，发起者会排出列表 \emph{by index in a loop} 而不是批量消耗向量——因为在排出运行时，允许一个回调对另一个回调进行排队。因此，执行顺序是明确定义的（FIFO 包括在排出期间排队的回调），并且排出在启动线程从 \texttt{backward} 返回之前运行，因此后向逻辑无需任何用户端轮询即可观察到完全累积的叶梯度。

\subsubsection{可重入向后和本地队列技巧}
\label{subsubsec:reentrant}

节点的 \texttt{apply} 本身可以调用 \texttt{backward}——与 \texttt{create\_graph} 的高阶微分，或减少图中损失的自定义函数。引擎跟踪 \texttt{thread\_local} 嵌套深度；当它非零时，执行路径将内部图的 \emph{every} 任务路由到仅当前线程耗尽的堆栈本地 \texttt{ReadyQueue} 。这样可以防止的死锁是结构性的：如果没有它，内部任务可能会落在设备队列上，而该队列的工作人员正忙于外部图，等待一个只能在内部向后返回后运行的节点——跨两个队列的循环。使用本地队列，内部图是严格单线程的、独立的，并且不能阻塞外部图工作线程；嵌套图会失去图内并行性，这是可以接受的，因为向后重入很少见且通常很小，而替代方案是调度优先级无法恢复的死锁。

\subsubsection{引擎跟踪：观察调度程序}
\label{subsubsec:enginetrace}

调度程序的决策——依赖计数、缓冲区合并、队列路由——正是读者无法从静态图中推断出的部分。环境门控跟踪设施使它们可观察：\texttt{TP\_ENGINE\_TRACE=1} 发出生命周期事件，level~2 添加每个节点应用程序的形状和每个交付决策（计数器递减、缓冲区命中、排队），level~3 打印截断的梯度 \emph{values}。每行都以图形的单调 id 为前缀，因此嵌套和并发向后在日志中保持可分离；节点被标记为 \texttt{name\#sequence\_nr@pointer-suffix}，该标签对于一个图保持稳定并且可在其上进行 grep 操作。当级别变量读取为零（每个事件站点一个静态整数加载）时，该设施编译为零成本，因此它可以在生产中启用并打开以应对生产中行为不当的图表，而无需重建。

\begin{lstlisting}[caption={Trace output is designed to be read: graph id, stable node labels, and the decision being made.}, style=pystyle]
# TP_ENGINE_TRACE=2 python train.py 2> trace.log
[tp-engine g1] begin execute create_graph=0 keep_graph=1 roots=1
[tp-engine g1] dep   MulBackward#38 -> AccumulateGrad#1 (count=1)
[tp-engine g1] buf   MulBackward#38 -> AccumulateGrad#1@7f2a (waiting on more inputs)
[tp-engine g1] apply MulBackward#38 [8,16] f32 cpu
[tp-engine g1] end   g1 outstanding=0
\end{lstlisting}

%----------------------------------------------------------------------

\subsection{二维显式 DAG}
\label{subsec:dagfig}

图~\ref{fig:autograddag}展示了引擎实际行走的图表。每个节点的\texttt{next\_edges\_}是传入的grad通道；堆键是 \texttt{sequence\_nr\_}。设备放置是一个 \texttt{Edge} 提示，而不是单独的计划。

\begin{figure}[H]
\centering
\begin{tikzpicture}[every node/.style={font=\scriptsize}, node distance=0.55cm, >=Stealth]
\tikzset{libbox/.style={draw, rounded corners=2pt, thick, align=center, font=\scriptsize, minimum height=0.75cm}}
\node[libbox, fill=blue!9, draw=tpblue, minimum width=3.2cm] (g1) {\texttt{GraphRoot} \tiny $(topo=5)$};
\node[libbox, fill=blue!9, draw=tpblue, minimum width=3.2cm, below=of g1] (n1) {\texttt{SumBackward} \tiny $seq=42$};
\node[libbox, fill=blue!9, draw=tpblue, minimum width=3.2cm, below=of n1] (n2) {\texttt{MulBackward} \tiny $seq=38$};
\node[libbox, fill=green!10, draw=tpgreen, below left=0.65cm and -0.3cm of n2, minimum width=2.9cm] (n3) {\texttt{AccumulateGrad(x)} \tiny $seq=1$};
\node[libbox, fill=purple!10, draw=tppurple, below right=0.65cm and -0.3cm of n2, minimum width=2.9cm] (n4) {\texttt{CatBackward} \tiny $seq=29$\\ \tiny \texttt{SavedVariable\{self,dim\}}};
\node[libbox, fill=gray!10, draw=gray, below=1.15cm of n3, minimum width=6.8cm] (q) {ReadyQueue CPU (max-heap by \texttt{sequence\_nr}) \quad ReadyQueue CUDA:0 (worker thread)};
\path[arrow] (g1) edge node[right, font=\tiny, color=gray] {\texttt{Edge\{hint (2,3)/fp32/cuda:0\}}} (n1);
\path[arrow] (n1) edge (n2);
\path[arrow] (n2) edge node[left, font=\tiny, color=gray] {\texttt{input\_nr 0}} (n3);
\path[arrow] (n2) edge node[right, font=\tiny, color=gray] {\texttt{input\_nr 1}} (n4);
\node[font=\tiny, color=gray, below=0.18cm of g1, xshift=2.35cm] {\texttt{dependencies: Mul 1, Sum 1}};
\node[font=\tiny, color=gray, below=0.02cm of q] {\texttt{tls nested\_depth\_} routes reentrant backs to \texttt{local\_queue}};
\draw[arrow, dashed, color=gray] (g1.east) -- ++(0.75,0) |- (q.east);
\draw[arrow, dashed, color=gray] (n1.east) -- ++(0.55,0) |- (q.east);
\draw[arrow, dashed, color=gray] (n2.east) -- ++(0.35,0) |- (q.east);
\end{tikzpicture}
\caption{显式 DAG，无磁带。节点携带 \texttt{ OutputSlotMeta}；边带有 \texttt{ device\_hint}。引擎的 \texttt{ReadyQueue} 首先弹出最大的 \texttt{sequence\_nr}； \texttt{dependencies\_} 大门准备就绪。可重入向后绕过设备队列进入由 \texttt{nested\_depth\_} 保护的堆栈本地 \texttt{local\_queue}。}
\label{fig:autograddag}
\end{figure}

%==============================================================================
\subsection{不变量、跟踪和本地验证}
\label{subsec:verify}

检查代码强制执行的不变量；违规行为要么被警告，要么被硬失败：

\begin{itemize}[leftmargin=1.2em]
\item 推送前更新。忘记这一点会让 \texttt{min\_topo} 错误地修剪并删除实时依赖项。
\item 需要 \texttt{use\_count==2} （临时 impl 别名算作 1）和 \texttt{storage==1}。在共享存储上就地累积会别名损坏同级。
\item 检查 \texttt{stop} \emph{after} 堆空测试，但 \emph{before} 等待。颠倒顺序会在图表完成后休眠一个额外的轮询间隔。
\item 必须先于队列推送。 Worker 逐步递减并测试 \texttt{outstanding==0}；不计数而推动会挂起启动器。
\item \texttt{saved\_version\_} 和当前版本都会抛出异常。该消息与用户在两个代码库中看到的字符串相同，从而启用传输的 Runbook。
\end{itemize}

关闭时跟踪是免费的（level~0：一次静态 \texttt{bool} 检查）。通过 \texttt{3} 启用：

\begin{lstlisting}[caption={Engine tracing (zero cost when off).}, label=lst:trace]
TP_ENGINE_TRACE=2 python -c "import tensorplay; x=tensorplay.randn([2,3], requires_grad=True); y=(x*2).sum(); y.backward()"
# [tp-engine g1] execute create_graph=0 keep_graph=0 accumulate_grad=1 roots=1
# [tp-engine g1] dep  SumBackward#42@suffix -> MulBackward#38@suffix (count=1)
# [tp-engine g1] apply MulBackward#38 inputs=[ones(2,3) fp32]  |  emit MulBackward#38 outputs=[ones(2,3)]
# [tp-engine g1] buf  MulBackward#38 -> AccumulateGrad#1@... (waiting on more inputs)
\end{lstlisting}

验证片段（全部传递引用的提交；结构检查不需要 Python 导入）：

\begin{lstlisting}[caption={Local reproduction of every cited path:line.}, label=lst:verify, style=pystyle]
grep -n "get_and_increment_sequence_nr\|sequence_nr_\|topological_nr_\|materialize_grads_\|output_metas_\|pre_hooks_\|anomaly_metadata_" tpx/include/Node.h | head -n 25
grep -n "shape_hint\|grad_dtype\|device_type_hint\|device_index_hint\|has_input_metadata" tpx/include/Edge.h
grep -n "not_ready_\|dependencies_\|exec_info_\|captured_vars_\|mutex_\|outstanding_tasks_" tpx/include/GraphTask.h
grep -n "can_accumulate_inplace\|void accumulate\|void add\|device_index" tpx/include/InputBuffer.h
grep -n "priority_queue<NodeTask>\|operator<.*sequence_nr\|pop_until\|wait_for.*5ms\|queue_for_device\|worker_main\|nested_depth_" tpx/include/Engine.h tpx/src/Engine.cpp | head -n 30
grep -n "compute_dependencies\|compute_min_topological_nr\|init_to_execute\|evaluate_function\|sum_to_shape\|tensor_has_nan\|GradModeGuard\|CurrentEvalNodeGuard\|release_variables" tpx/src/Engine.cpp | head -n 40
grep -n "class SavedVariable\|save(\|unpack()\|saved_version_\|reset_data" tpx/include/SavedVariable.h tpx/src/SavedVariable.cpp
grep -n "VersionCounter\|VariableVersion\|shares_counter\|share_version_counter\|bump_version" p10/include/VariableVersion.h p10/include/TensorImpl.h | head -n 30
grep -n "AnomalyMode\|AnomalyMetadata\|tls_current_evaluating_node\|CurrentEvalNodeGuard\|store_stack\|print_stack" tpx/include/AnomalyMode.h tpx/src/AnomalyMode.cpp | head -n 30
grep -n "GradMode::is_enabled\|InferenceMode::is_enabled\|thread_local.*enabled_" p10/include/GradMode.h p10/include/InferenceMode.h p10/src/GradMode.cpp p10/src/InferenceMode.cpp | head
grep -n "class PyNode\|num_inputs.*output_metas\|apply.*variable_list\|attach_outputs\|custom_function_apply\|run_custom_function_forward" src/bindings/python/Autograd.cpp | head -n 40
# quick functional smoke (engine + version guard + anomaly)
python3 - <<'PY'
import tensorplay as tp
x = tp.randn([4], requires_grad=True)
y = tp.randn([4], requires_grad=True)
z = (x * y).sum(); z.backward()
assert x.grad is not None and y.grad is not None
# in-place guard fires
try:
    a = tp.randn([3], requires_grad=True); a_saved = a
    b = (a * 2).sum(); a.add_(1); b.backward()
    raise SystemExit("expected version error")
except RuntimeError as e:
    assert "modified by an inplace operation" in str(e)
# anomaly includes forward traceback
tp.autograd.set_detect_anomaly(True)
try:
    p = tp.randn([2], requires_grad=True); q = (p * p).sum(); q.backward()
except Exception: pass
finally: tp.autograd.set_detect_anomaly(False)
print("autograd invariants hold")
PY
\end{lstlisting}

引擎的 \texttt{ReadyQueue} 启发式加上精确的 \texttt{dependencies\_} 门控共同保证在其消费者全部交付之前没有节点运行，而远离根的节点（大 \texttt{sequence\_nr}）会尽早运行，限制 \texttt{not\_ready\_} 占用。版本防护保证就地突变后不会出现静默错误的等级。全局异常模式加上 \texttt{tls\_current\_evaluating\_node} 保证向后创建的子图中的失败可归因于引发它的前向 Python 调用站点。这三种机制（显式 DAG、版本化存储、全局异常链）共同构成了框架其余部分所依赖的契约。